\documentclass[11pt]{article}
\usepackage[margin=1in]{geometry}
\usepackage{amsmath,amssymb}
\usepackage{hyperref}
\usepackage{booktabs}

\title{$AnAn$ versus PL$(AnAn)$ photon fluxes:\\
why our plot and arXiv:2606.05469 look different}
\author{}
\date{\today}

\begin{document}
\maketitle

\section{Summary}

In \texttt{plotting\_scripts/flux\_comparison.py} our $AnAn$ flux crosses \emph{above}
the point-like PL$(AnAn)$ flux at $z\gtrsim0.03$, whereas in
Fig.~\texttt{flux\_AA\_classes} of Guzey, Innocenti, Sta\'sto and Strikman
(arXiv:2606.05469) the PL flux is \emph{always} the larger one. The two plots
do not contradict each other: the curves labelled $AnAn$ are different
quantities.
\begin{itemize}
  \item Their $AnAn$ flux counts every photon at the centre of the target
    nucleus (single impact parameter $b$, weighted by $\Gamma_{AA}(b)$).
  \item Our $AnAn$ curve is the \emph{effective flux} of Eskola et al.\
    (arXiv:2404.09731, Eq.~4), which averages over the positions of the
    target nucleons.
\end{itemize}
Evaluated with the same single-$b$ definition as theirs, our code also gives
a flux below PL$(AnAn)$ at every $z$.

\section{Why the smallest distance controls large $z$}

For the point-like charge distribution,
\begin{equation}
  f^{\rm PL}_{\gamma/A}(z,b) = \frac{\alpha_{\rm em}Z^2}{\pi^2}\,\frac1z
  \left(z m_N K_1(z m_N b)\right)^2 ,
\end{equation}
and for large arguments $K_1(x)\sim\sqrt{\pi/2x}\,e^{-x}$, so
\begin{equation}
  f^{\rm PL}_{\gamma/A}(z,b) \propto e^{-2 z m_N b}.
\end{equation}
Physically, the Lorentz-contracted field of the nucleus is a pulse, and a
pulse passing at distance $b$ only contains frequencies up to
$\omega\sim\gamma/b$. High-energy photons (large $z$) therefore come only
from the closest encounters, and the $b$-integrated flux at large $z$ is
exponentially sensitive to the smallest distance that is allowed.

\section{The three fluxes}

\paragraph{PL$(AnAn)$.} Point-like flux, measured at the target centre,
with a sharp cut at $b_{\min}=14.2$~fm (\texttt{photon\_flux()} in
\texttt{src/photon\_flux.cpp}; Eq.~(\ref{eq:flux_PL}) of the paper):
\begin{equation}
  f^{{\rm PL}(AnAn)}(z) = \int_{b\ge b_{\min}} d^2b\; f^{\rm PL}_{\gamma/A}(z,b).
  \label{eq:flux_PL}
\end{equation}

\paragraph{$AnAn$, single $b$ (the paper's definition).}
\begin{equation}
  f^{(AnAn)}(z) = \int d^2b\; f_{\gamma/A}(z,b)\,\Gamma_{AA}(b),
\end{equation}
where $b$ is ``the transverse distance between the centers of the two
nuclei'' (paper, Sec.~3). The Woods--Saxon form factor in $f_{\gamma/A}$
changes nothing here: for $b>R_A$ the Woods--Saxon flux equals the
point-like one (the field outside a spherically symmetric charge does not
depend on its internal structure), and $\Gamma_{AA}$ removes everything
inside that range anyway.

\paragraph{$AnAn$, effective flux (our plot, our default \texttt{FLUX\_MODEL=EFF}).}
\begin{equation}
  f^{(AnAn)}_{\rm eff}(z) = \frac{1}{B}\int d^2r\; f_{\gamma/A}(z,r)
  \int d^2s\; T_B(s)\,\Gamma_{AA}(|\mathbf r-\mathbf s|),
\end{equation}
where $\mathbf r$ is the photon's position relative to the emitter,
$\mathbf s$ the position of a target nucleon relative to the target centre
and $b=|\mathbf r-\mathbf s|$ the centre-to-centre distance. Nucleons on the
near side of the target sit a few fm closer to the emitter than the target
centre does.

\paragraph{The paper's own comment.} The paper cites exactly this
refinement and chooses not to use it: ``One can in principle refine
[their single-$b$ $AnAn$ formula] by taking into account details of the
transverse-plane geometry of UPC events \dots\ While it somewhat increases
$f_{\gamma/A}(z)$ for large $z$, the photon flux is suppressed by several
orders of magnitude in this region.''

\section{Why their PL flux is always larger}

Both of their curves measure the photon at the target centre; they differ
only in the cut at small $b$:
\begin{center}
\begin{tabular}{lcccc}
\toprule
$b$ [fm] & $\le 13.8$ & 14.2 & 15.5 & $\gtrsim 19$ \\
\midrule
sharp cut (PL) & 0 & 1 & 1 & 1 \\
$\Gamma_{AA}(b)$ (\texttt{input/Gamma\_AA.dat}) & 0 & $\approx0.01$ & $\approx0.5$ & $\approx1$ \\
\bottomrule
\end{tabular}
\end{center}
$\Gamma_{AA}$ removes strictly more than the sharp cut, so the single-$b$
$AnAn$ flux is below PL$(AnAn)$ at every $z$. The gap grows with $z$ because
large-$z$ photons come from the smallest $b$, which is exactly where
$\Gamma_{AA}$ suppresses.

\section{Why our effective flux crosses above PL at large $z$}

\subsection{One exponential, evaluated at two different distances}

The $e^{-2zm_N b}$ of Sec.~2 and the $e^{-2zm_N r}$ used below are the
\emph{same} function. Writing $\kappa \equiv 2 z m_N$,
\begin{equation}
  f_{\gamma/A}(z,d) \propto e^{-\kappa d},
  \label{eq:exp}
\end{equation}
where $d$ is the distance from the centre of the emitting nucleus to the
point where the photon is absorbed. The two treatments only differ in what
they use for $d$:
\begin{itemize}
  \item single $b$ (PL$(AnAn)$ and the paper's $AnAn$): the photon is
    absorbed at the target centre, $d=b$;
  \item effective flux: the photon is absorbed by a nucleon at position
    $\mathbf s$ inside the target, $d=r=|\mathbf b+\mathbf s|$.
\end{itemize}
Equation~(\ref{eq:exp}) uses the large-argument form of $K_1$, valid for
$z m_N d \gtrsim 1$; with $d\approx14$~fm that means $z\gtrsim0.015$. The
effective lengths quoted below are only meaningful in that range.

\subsection{The loss from $\Gamma_{AA}$: a fixed length}

Replacing the sharp cut at 14.2~fm by $\Gamma_{AA}(b)$ (which rises from 0
at 14~fm to 0.5 at 15.5~fm and 1 at $\sim19$~fm) effectively moves the
closest allowed distance outward by some $\delta_{\rm loss}$:
\begin{equation}
  L(z) \equiv \frac{f^{(AnAn)}_{\rm single\ b}}{f^{{\rm PL}(AnAn)}}
  \approx e^{-\kappa\,\delta_{\rm loss}}, \qquad
  \delta_{\rm loss} = -\frac{\ln L}{\kappa}.
\end{equation}
$\delta_{\rm loss}$ is set by the width of the $\Gamma_{AA}$ edge and is
roughly independent of $z$: about 1.2--1.5~fm (table below).

\subsection{The gain from averaging over nucleons: a growing length}

For $b\gg|\mathbf s|$, a nucleon at longitudinal offset $s_\parallel$
towards the emitter sits at $r\approx b-s_\parallel$, so its flux is
enhanced by $e^{\kappa s_\parallel}$ relative to the target centre.
Averaging over the target thickness $T_B$,
\begin{equation}
  G(z) \equiv \frac{f^{(AnAn)}_{\rm eff}}{f^{(AnAn)}_{\rm single\ b}}
  \approx \left\langle e^{\kappa s_\parallel} \right\rangle_{T_B}
  \equiv e^{\kappa\,\delta_{\rm gain}(z)} .
\end{equation}
Two properties follow:
\begin{itemize}
  \item $G\ge1$ always. The target is symmetric, $\langle s_\parallel\rangle=0$,
    but the exponential is convex, so the near-side nucleons gain more than
    the far-side ones lose: $\langle e^{\kappa s_\parallel}\rangle \ge
    e^{\kappa\langle s_\parallel\rangle}=1$ (Jensen's inequality).
  \item $\delta_{\rm gain}$ grows with $z$. For small $\kappa$,
    $\delta_{\rm gain}\approx\kappa\langle s_\parallel^2\rangle/2\to0$:
    long-range photons see the whole target and the offsets average out.
    For large $\kappa$ the average is dominated by the nucleons closest to
    the emitter, and $\delta_{\rm gain}$ tends to the depth of the target's
    near edge, a few fm.
\end{itemize}

\subsection{Putting them together}

The ratio plotted in \texttt{flux\_comparison.py} is the product of the two:
\begin{equation}
  \frac{f^{(AnAn)}_{\rm eff}}{f^{{\rm PL}(AnAn)}} = L(z)\,G(z)
  \approx e^{\kappa\,(\delta_{\rm gain}(z)-\delta_{\rm loss})} .
\end{equation}
Both factors are exponential in $z$ with the same $\kappa$; what differs is
the length multiplying it. $\delta_{\rm loss}$ is fixed, $\delta_{\rm gain}$
grows. The effective flux therefore sits below PL while
$\delta_{\rm gain}<\delta_{\rm loss}$ and crosses above it once
$\delta_{\rm gain}$ overtakes, at $z\approx0.02$.

The single-$b$ scan isolates the two factors cleanly: it has the same
$\Gamma_{AA}$ as the effective flux but no nucleon average. From our scans:
\begin{center}
\begin{tabular}{ccccccc}
\toprule
$z$ & $L$ & $G$ & $L\,G$ & $\delta_{\rm loss}$ [fm] & $\delta_{\rm gain}$ [fm] & $\delta_{\rm gain}-\delta_{\rm loss}$ [fm] \\
\midrule
0.001 & 0.958 & 1.019 & 0.977 & -- & -- & -- \\
0.01  & 0.848 & 1.115 & 0.945 & -- & -- & -- \\
0.02  & 0.747 & 1.323 & 0.988 & 1.53 & 1.47 & $-0.06$ \\
0.03  & 0.660 & 1.690 & 1.115 & 1.46 & 1.84 & 0.38 \\
0.05  & 0.522 & 3.390 & 1.770 & 1.37 & 2.57 & 1.20 \\
0.07  & 0.420 & 8.196 & 3.442 & 1.30 & 3.16 & 1.86 \\
0.1   & 0.312 & 45.07 & 14.04 & 1.23 & 4.01 & 2.78 \\
\bottomrule
\end{tabular}
\end{center}
(Lengths are omitted for $z<0.015$, where Eq.~(\ref{eq:exp}) does not apply.)
The loss factor alone reproduces the paper's picture (PL always on top).
The crossover in our plot comes entirely from $G$, whose effective length
grows from $\sim1.5$~fm to $\sim4$~fm over this range.

\section{Numbers from our code}

Values of $dN/dz$ from \texttt{out/flux\_scan.csv} (PL$(AnAn)$ and single-$b$
$AnAn$) and \texttt{out/flux\_scan\_eff\_AnAn.csv} (effective $AnAn$), Pb--Pb
at $\sqrt{s_{NN}}=5.36$~TeV. Values between scan grid points are obtained by
interpolating each flux's ratio to the analytic PL formula, not the flux
itself (log-log interpolation of the steeply falling flux on the coarse grid
underestimates it at large $z$):
\begin{center}
\begin{tabular}{ccccc}
\toprule
$z$ & PL$(AnAn)$ & $AnAn$, single $b$ & $AnAn$, effective & effective/PL \\
    &            & (= paper)           & (= our plot)      & \\
\midrule
0.001 & $7.25\times10^{4}$ & $6.95\times10^{4}$ & $7.08\times10^{4}$ & 0.98 \\
0.003 & $1.32\times10^{4}$ & $1.23\times10^{4}$ & $1.27\times10^{4}$ & 0.96 \\
0.01  & $1.07\times10^{3}$ & $9.07\times10^{2}$ & $1.01\times10^{3}$ & 0.95 \\
0.02  & $1.16\times10^{2}$ & $8.68\times10^{1}$ & $1.15\times10^{2}$ & 0.99 \\
0.03  & $18.5$             & $12.2$             & $20.6$             & 1.12 \\
0.05  & $0.685$            & $0.358$            & $1.21$             & 1.77 \\
0.07  & $3.15\times10^{-2}$ & $1.32\times10^{-2}$ & $0.108$           & 3.4 \\
0.1   & $3.71\times10^{-4}$ & $1.16\times10^{-4}$ & $5.21\times10^{-3}$ & 14.0 \\
\bottomrule
\end{tabular}
\end{center}
\begin{itemize}
  \item The single-$b$ column is below PL$(AnAn)$ everywhere, reproducing
    the paper's figure: close to PL for $z<0.01$, separating at large $z$.
  \item The effective flux is a few percent below PL for $z\lesssim0.02$
    (the stricter $\Gamma_{AA}$ wins; long-range photons do not resolve the
    target's size) and grows exponentially above it. A factor 14 at
    $z=0.1$ corresponds to an effective closest distance about 2.8~fm
    shorter, from $\ln 14 \approx 2 z m_N \Delta b$ (Sec.~5).
\end{itemize}

\section{Consequences for our $D^0$ calculation}

\begin{itemize}
  \item The smallest photon fraction probed at a given $(p_T,y)$ is
    $z_{\min}=m_T e^{y}/\sqrt{s_{NN}}$: for $p_T=2$~GeV, $z_{\min}\approx0.004$
    at $y=2$ and $\approx0.03$ at $y=4$. Most of our kinematics lie where the
    fluxes agree to a few percent; a Pb--Pb test point at $p_T=2$~GeV, $y=0$
    gave a $\sim2\%$ difference between \texttt{EFF} and \texttt{PL}.
  \item The default \texttt{D0} setting (\texttt{FLUX\_MODEL=EFF}) is
    \emph{not} the paper's flux. To reproduce their treatment, run with
    \texttt{FLUX\_MODEL=WS} (identical to \texttt{PL} in Pb--Pb), which
    computes $\int d^2b\, f_{\gamma/A}(z,b)\,\Gamma_{AA}(b)\,P_{\text{no-EM}}(b)$,
    their $An0n$ formula.
  \item Differences between the two treatments show up only at large $z$,
    i.e.\ forward rapidity and high $p_T$.
\end{itemize}

\end{document}
